%% Chapitre 03 — « Ta première tâche réelle déléguée »
%% RÈGLE : uniquement des commandes sémantiques bookforge. Aucun \chapter,
%% \textbf, \begin{lstlisting}, \vspace... brut.

\chapitre{Ta première tâche réelle déléguée}

\introChapitre{
  Fini les essais à blanc : on confie à l'agent une vraie petite tâche sur du
  vrai code. L'objectif n'est pas le résultat parfait mais d'observer comment
  l'agent explore, propose, modifie — et de prendre le réflexe de relire avant
  d'accepter. C'est ici que naît la confiance mesurée.
}

\section{Choisir une bonne première tâche}

\paragraphe{Au chapitre précédent, tu as installé Claude Code, lancé une
première session et repéré l'interface. Tout cela restait à blanc : on
regardait l'agent répondre, sans rien lui faire toucher. Place maintenant à du
concret. On va lui confier une vraie modification sur un vrai dépôt, et tu vas
décider, ligne par ligne, ce qui entre dans ton code.}

\paragraphe{La première tâche que tu délègues n'a pas vocation à t'impressionner.
Elle a vocation à t'apprendre à lire l'agent. Choisis donc quelque chose de
petit, de vérifiable et de peu risqué. Petit : quelques fichiers au plus, pas
une refonte. Vérifiable : tu sais dire si c'est juste, idéalement en lançant un
test ou en relisant à l'œil. Peu risqué : si ça part en vrille, tu jettes le
travail sans avoir cassé ta journée.}

\paragraphe{Quelques candidates idéales pour un premier vol :}

\begin{liste}
  \element{Ajouter une fonction utilitaire isolée, avec son test.}
  \element{Corriger un message d'erreur trompeur ou une faute dans la
  documentation.}
  \element{Extraire une portion de code dupliquée dans une fonction partagée.}
  \element{Couvrir d'un test un cas limite oublié.}
\end{liste}

\paragraphe{Ce qu'il faut éviter pour un premier essai : tout ce qui touche à
la sécurité, à la facturation, aux migrations de base de données, ou qui
s'étale sur dix fichiers couplés. Ces tâches sont déléguables aussi, mais plus
tard, quand tu auras le réflexe de relecture bien ancré.}

\begin{astuce}
Travaille sur une branche dédiée, ou au minimum sur un dépôt dont l'état est
propre (\code{git status} ne montre rien en attente). Ainsi, quoi que propose
l'agent, un \code{git diff} ou un \code{git checkout .} te ramène à zéro en une
commande. C'est ton filet de sécurité le plus simple, et le plus efficace.
\end{astuce}

\section{Le dépôt-exemple : formuler une intention, pas des instructions}

\paragraphe{Prenons un dépôt concret pour la suite du chapitre. Imagine une
petite API en Python, \code{ledger}, qui tient des comptes : un module
\code{accounts.py} avec une fonction \code{transfer(src, dst, amount)} qui
déplace un montant d'un compte à l'autre. Un détail t'agace : quand le montant
est négatif, la fonction l'accepte sans broncher et crédite le mauvais compte.
Tu veux qu'un montant négatif ou nul soit refusé proprement.}

\paragraphe{Le réflexe du débutant, c'est de dicter la solution : « ajoute un
\code{if amount <= 0} au début de \code{transfer} et lève un
\code{ValueError} ». Ça marche, mais tu te prives de l'agent. Tu lui imposes ta
première idée alors qu'il pourrait repérer que la même validation manque
ailleurs, ou qu'il existe déjà une exception maison pour ce genre de cas.}

\paragraphe{Formule plutôt une \important{intention} : le but, et le critère de
réussite. Laisse l'agent décider du « comment » après avoir lu le code. Une
demande bien calibrée ressemble à ceci :}

\begin{extraitcode}{text}
Dans accounts.py, transfer() accepte des montants négatifs ou nuls et
produit des transferts incohérents. Refuse ces montants avec une erreur
claire. Regarde d'abord comment le reste du module signale ses erreurs
pour rester cohérent, et ajoute un test qui couvre le cas du montant nul
et du montant négatif.
\end{extraitcode}

\paragraphe{Remarque ce qu'on a donné : une intention nette (refuser les
montants invalides), un critère de réussite (un test qui passe), et une
consigne de cohérence (s'aligner sur l'existant). On n'a pas dicté le nom de
l'exception ni l'emplacement exact du \code{if}. C'est l'agent qui va le
déterminer en lisant — et c'est exactement ce qu'on veut observer.}

\section{Observer l'agent explorer et proposer un plan}

\paragraphe{Une fois la demande envoyée, l'agent ne se jette pas sur le clavier.
Comme on l'a vu au premier chapitre, il agit en boucle : il lit, il raisonne,
il agit. Sa première action ici est presque toujours d'explorer le dépôt pour
se constituer un contexte. Tu le vois lire des fichiers, parfois lancer une
recherche dans le code, avant même d'écrire la moindre ligne.}

\paragraphe{Dans notre exemple, la session ressemble à ceci :}

\begin{extraitcode}{text}
> Je regarde accounts.py pour comprendre transfer() et le style d'erreurs.
  [lecture de accounts.py]
  [lecture de errors.py]
  [recherche « raise » dans le module]

Le module définit déjà une exception InvalidAmount dans errors.py,
utilisée par withdraw(). transfer() ne fait aucune validation. Je vais :
  1. valider amount > 0 au début de transfer(), en levant InvalidAmount ;
  2. ajouter deux tests dans test_accounts.py (montant nul, négatif).
Je te montre le changement avant de l'appliquer.
\end{extraitcode}

\paragraphe{Cet instant vaut de l'or. L'agent a trouvé tout seul l'exception
\code{InvalidAmount} que tu avais peut-être oubliée, et il s'aligne dessus au
lieu d'inventer un \code{ValueError} générique. C'est le bénéfice direct de
l'intention plutôt que de la micro-instruction : il a comblé un manque de ton
contexte à toi.}

\begin{note}
Lis toujours le plan annoncé avant le code. Un plan qui parle de fichiers que
tu ne t'attendais pas à voir touchés est un signal d'alarme précoce : tu peux
recadrer immédiatement, avant qu'une seule ligne ne soit écrite. C'est bien
moins coûteux que de défaire un changement déjà appliqué.
\end{note}

\section{Lire un diff et décider}

\paragraphe{L'agent passe ensuite à l'action et te présente ses modifications
sous forme de \terme{diff} : une vue ligne par ligne de ce qu'il ajoute et de
ce qu'il retire. Les lignes précédées d'un \code{+} sont ajoutées, celles
précédées d'un \code{-} sont supprimées ; le reste est du contexte inchangé,
affiché pour que tu te repères. C'est l'unité de base que tu vas relire et
approuver pendant tout le reste de ce livre, alors prenons le temps de bien la
lire.}

\begin{extraitcode}{text}
  accounts.py
  --- a/accounts.py
  +++ b/accounts.py
@@ def transfer(src, dst, amount):
-     def transfer(src, dst, amount):
-         src.balance -= amount
-         dst.balance += amount
+     def transfer(src, dst, amount):
+         if amount <= 0:
+             raise InvalidAmount(amount)
+         src.balance -= amount
+         dst.balance += amount
\end{extraitcode}

\paragraphe{Que regardes-tu dans un diff ? Trois choses, dans l'ordre. D'abord,
est-ce que ça fait bien ce qui était demandé — ici, refuser les montants
invalides ? Oui. Ensuite, est-ce que ça ne fait \important{rien d'autre} en
douce — un fichier de configuration modifié au passage, une dépendance
ajoutée ? Non, le changement est chirurgical. Enfin, est-ce cohérent avec le
projet — l'exception est bien celle du module, l'importation existe déjà ? Oui.}

\paragraphe{Le second diff, celui du test, mérite la même attention. Un agent
peut écrire un test qui passe pour de mauvaises raisons, par exemple un test
qui n'exécute jamais l'assertion. Relis-le comme du vrai code :}

\begin{extraitcode}{text}
  test_accounts.py
+ def test_transfer_montant_negatif_refuse():
+     a, b = Account(100), Account(0)
+     with pytest.raises(InvalidAmount):
+         transfer(a, b, -10)
+
+ def test_transfer_montant_nul_refuse():
+     a, b = Account(100), Account(0)
+     with pytest.raises(InvalidAmount):
+         transfer(a, b, 0)
\end{extraitcode}

\paragraphe{Là, tu as trois décisions possibles, et une seule à prendre :}

\begin{liste}
  \element{\important{Accepter} : le diff fait le travail, proprement. Tu
  approuves, l'agent applique, tu lances les tests pour confirmer.}
  \element{\important{Corriger} : le fond est bon mais un détail cloche (un nom
  maladroit, un cas oublié). Tu le dis en une phrase et l'agent ajuste.}
  \element{\important{Refuser} : la direction est mauvaise. Tu refuses, tu
  réexpliques l'intention, et l'agent repart sur une autre base.}
\end{liste}

\begin{avertissement}
Ne prends jamais l'habitude d'accepter un diff sans le lire parce qu'il « a
l'air bien » ou que les tests passent. Un test écrit par l'agent et un code
écrit par l'agent peuvent être faux du même biais et se valider mutuellement.
La relecture du diff par un humain reste ton dernier rempart. Le jour où tu
cliques « accepter » par réflexe, tu as troqué ta confiance mesurée contre de
la confiance aveugle.
\end{avertissement}

\section{Le cycle court : demander, observer, ajuster}

\paragraphe{Tu viens de vivre une boucle complète : tu as demandé, l'agent a
exploré et proposé, tu as relu et décidé. C'est le rythme de base du travail
avec un agent, et il a une vertu : il est \important{court}. Tu n'attends pas
qu'un gros chantier soit fini pour découvrir qu'il a déraillé dès la troisième
ligne.}

\paragraphe{Supposons que le diff t'ait presque convenu, mais que tu veuilles
un message d'erreur plus parlant. Pas besoin de tout reprendre. Tu réponds dans
la même session, qui a gardé tout son contexte :}

\begin{extraitcode}{text}
> Bien. Mais InvalidAmount(amount) donne un message obscur en log.
  Passe-lui un texte explicite, du genre « montant invalide : doit être
  strictement positif ». Garde le reste tel quel.
\end{extraitcode}

\paragraphe{L'agent ne recommence pas de zéro : il ajuste le diff existant et te
le re-présente. Cette conversation incrémentale, c'est tout l'intérêt de rester
dans une session continue plutôt que de relancer une demande géante à chaque
correction. Plus tes boucles sont courtes, plus tu gardes la main.}

\begin{astuce}
Tant que tu apprends à lire un agent, garde tes demandes à la taille d'un seul
diff relisible d'un coup d'œil. Si une tâche produit un diff trop gros pour être
relu sans fatigue, c'est le signe qu'il fallait la découper. On en fera une
compétence à part entière au chapitre suivant.
\end{astuce}

\section{Débrief : ce qui a marché, ce qui a dérapé}

\paragraphe{Prends l'habitude, après chaque tâche déléguée, d'un court débrief
mental. Pas de cérémonie : deux questions suffisent. Qu'est-ce que l'agent a
fait mieux ou plus vite que toi ? Et où a-t-il failli te faire accepter quelque
chose de bancal ?}

\paragraphe{Sur notre exemple, le bilan est honnête. Au crédit de l'agent : il a
exploré seul, retrouvé l'exception maison que tu avais oubliée, et produit un
changement propre accompagné de ses tests. Tu as gagné du temps réel sur une
corvée réelle.}

\paragraphe{Mais imaginons un dérapage plausible, de ceux qu'on assume dans ce
livre. Dans une variante de la session, l'agent avait d'abord proposé de
modifier aussi \code{withdraw()} « pour homogénéiser », alors que personne ne
le lui avait demandé. Le diff était plus large que la tâche. Inoffensif ?
Peut-être. Mais c'est exactement le genre de débordement silencieux que la
relecture attrape : tu as recadré d'une phrase, et le périmètre est redevenu
celui que tu voulais. Sans relecture, ce surplus serait entré dans ton code
sans que tu le saches.}

\resumeChapitre{
  \element{\important{Intention, pas instruction} : formule le but et le critère de réussite, laisse l'agent décider du « comment » après avoir lu le code.}
  \element{\important{Plan avant le code} : lis toujours le plan annoncé par l'agent avant qu'il agisse — un périmètre inattendu se corrige en une phrase, un diff déjà appliqué se défait bien plus difficilement.}
  \element{\important{Lire un diff} : vérifie que ça fait ce qui était demandé, qu'il ne modifie rien d'autre en douce, et que c'est cohérent avec le projet.}
  \element{\important{Boucle courte} : reste dans la même session pour ajuster ; l'agent garde son contexte et affine sans repartir de zéro.}
  \element{\important{Débrief systématique} : après chaque tâche, note ce que l'agent a bien fait et où il a failli déborder — ce réflexe ancre l'habitude de relecture.}
}

\paragraphe{Retiens donc de ce premier vol moins le résultat que le réflexe :
intention claire en entrée, relecture du diff en sortie, boucle courte au
milieu. C'est ce trépied qui rend la délégation sûre, et tu vas le rejouer à
chaque chapitre. Reste une question : comment reconnaître, parmi tout ce que tu
fais dans une journée, ce qui se prête à ce traitement ? Comment penser ton
travail en morceaux qu'un agent peut réussir ? C'est précisément ce qu'on va
généraliser au chapitre suivant, en apprenant à penser en tâches que l'on peut
déléguer.}
